\chapter{Objetivos}
\label{pre:objetivos}

O objetivo deste trabalho \'e investigar e desenvolver um modelo de redes neurais de grafos capaz de melhorar o reconhecimento de padr\~oes em dados educacionais, com foco na identifica\c{c}\~ao precoce de problemas envolvendo alunos em risco de evas\~ao escolar e na identifica\c{c}\~ao de fatores que contribuem para o fracasso acad\^emico. Em perspectiva, a proposta visa ajudar institui\c{c}\~oes escolares a antecipar esses acontecimentos e construir m\'etodos para evit\'a-los.

As quest\~oes de pesquisa que orientam o estudo s\~ao:

\begin{enumerate}
    \item Como as redes neurais para grafos podem ser aplicadas com sucesso \`a an\'alise de dados educacionais?
    \item Em compara\c{c}\~ao com m\'etodos tradicionais, em que medida essa abordagem melhora a precis\~ao e a efici\^encia do reconhecimento de padr\~oes?
    \item Quais s\~ao as implica\c{c}\~oes pr\'aticas da aplica\c{c}\~ao de redes neurais gr\'aficas em ambientes educacionais?
    \item Quais s\~ao as limita\c{c}\~oes e os desafios dessa abordagem e como podem ser superados?
\end{enumerate}

Essas quest\~oes orientam a investiga\c{c}\~ao proposta e buscam contribuir para a melhoria da qualidade da educa\c{c}\~ao e para o enfrentamento de desafios institucionais recorrentes.
